% Options for packages loaded elsewhere
\PassOptionsToPackage{unicode}{hyperref}
\PassOptionsToPackage{hyphens}{url}
\PassOptionsToPackage{dvipsnames,svgnames,x11names}{xcolor}
\documentclass[
  11pt,
]{article}
\usepackage{xcolor}
\usepackage[margin=2cm]{geometry}
\usepackage{amsmath,amssymb}
\setcounter{secnumdepth}{-\maxdimen} % remove section numbering
\usepackage{iftex}
\ifPDFTeX
  \usepackage[T1]{fontenc}
  \usepackage[utf8]{inputenc}
  \usepackage{textcomp} % provide euro and other symbols
\else % if luatex or xetex
  \usepackage{unicode-math} % this also loads fontspec
  \defaultfontfeatures{Scale=MatchLowercase}
  \defaultfontfeatures[\rmfamily]{Ligatures=TeX,Scale=1}
\fi
\usepackage{lmodern}
\ifPDFTeX\else
  % xetex/luatex font selection
    \setmonofont[Scale=0.8,BoldFont=DejaVuSansMono-Bold.ttf]{DejaVuSansMono.ttf}
\fi
% Use upquote if available, for straight quotes in verbatim environments
\IfFileExists{upquote.sty}{\usepackage{upquote}}{}
\IfFileExists{microtype.sty}{% use microtype if available
  \usepackage[]{microtype}
  \UseMicrotypeSet[protrusion]{basicmath} % disable protrusion for tt fonts
}{}
\makeatletter
\@ifundefined{KOMAClassName}{% if non-KOMA class
  \IfFileExists{parskip.sty}{%
    \usepackage{parskip}
  }{% else
    \setlength{\parindent}{0pt}
    \setlength{\parskip}{6pt plus 2pt minus 1pt}}
}{% if KOMA class
  \KOMAoptions{parskip=half}}
\makeatother
\usepackage{color}
\usepackage{fancyvrb}
\newcommand{\VerbBar}{|}
\newcommand{\VERB}{\Verb[commandchars=\\\{\}]}
\DefineVerbatimEnvironment{Highlighting}{Verbatim}{commandchars=\\\{\}}
% Add ',fontsize=\small' for more characters per line
\newenvironment{Shaded}{}{}
\newcommand{\AlertTok}[1]{\textcolor[rgb]{1.00,0.00,0.00}{\textbf{#1}}}
\newcommand{\AnnotationTok}[1]{\textcolor[rgb]{0.38,0.63,0.69}{\textbf{\textit{#1}}}}
\newcommand{\AttributeTok}[1]{\textcolor[rgb]{0.49,0.56,0.16}{#1}}
\newcommand{\BaseNTok}[1]{\textcolor[rgb]{0.25,0.63,0.44}{#1}}
\newcommand{\BuiltInTok}[1]{\textcolor[rgb]{0.00,0.50,0.00}{#1}}
\newcommand{\CharTok}[1]{\textcolor[rgb]{0.25,0.44,0.63}{#1}}
\newcommand{\CommentTok}[1]{\textcolor[rgb]{0.38,0.63,0.69}{\textit{#1}}}
\newcommand{\CommentVarTok}[1]{\textcolor[rgb]{0.38,0.63,0.69}{\textbf{\textit{#1}}}}
\newcommand{\ConstantTok}[1]{\textcolor[rgb]{0.53,0.00,0.00}{#1}}
\newcommand{\ControlFlowTok}[1]{\textcolor[rgb]{0.00,0.44,0.13}{\textbf{#1}}}
\newcommand{\DataTypeTok}[1]{\textcolor[rgb]{0.56,0.13,0.00}{#1}}
\newcommand{\DecValTok}[1]{\textcolor[rgb]{0.25,0.63,0.44}{#1}}
\newcommand{\DocumentationTok}[1]{\textcolor[rgb]{0.73,0.13,0.13}{\textit{#1}}}
\newcommand{\ErrorTok}[1]{\textcolor[rgb]{1.00,0.00,0.00}{\textbf{#1}}}
\newcommand{\ExtensionTok}[1]{#1}
\newcommand{\FloatTok}[1]{\textcolor[rgb]{0.25,0.63,0.44}{#1}}
\newcommand{\FunctionTok}[1]{\textcolor[rgb]{0.02,0.16,0.49}{#1}}
\newcommand{\ImportTok}[1]{\textcolor[rgb]{0.00,0.50,0.00}{\textbf{#1}}}
\newcommand{\InformationTok}[1]{\textcolor[rgb]{0.38,0.63,0.69}{\textbf{\textit{#1}}}}
\newcommand{\KeywordTok}[1]{\textcolor[rgb]{0.00,0.44,0.13}{\textbf{#1}}}
\newcommand{\NormalTok}[1]{#1}
\newcommand{\OperatorTok}[1]{\textcolor[rgb]{0.40,0.40,0.40}{#1}}
\newcommand{\OtherTok}[1]{\textcolor[rgb]{0.00,0.44,0.13}{#1}}
\newcommand{\PreprocessorTok}[1]{\textcolor[rgb]{0.74,0.48,0.00}{#1}}
\newcommand{\RegionMarkerTok}[1]{#1}
\newcommand{\SpecialCharTok}[1]{\textcolor[rgb]{0.25,0.44,0.63}{#1}}
\newcommand{\SpecialStringTok}[1]{\textcolor[rgb]{0.73,0.40,0.53}{#1}}
\newcommand{\StringTok}[1]{\textcolor[rgb]{0.25,0.44,0.63}{#1}}
\newcommand{\VariableTok}[1]{\textcolor[rgb]{0.10,0.09,0.49}{#1}}
\newcommand{\VerbatimStringTok}[1]{\textcolor[rgb]{0.25,0.44,0.63}{#1}}
\newcommand{\WarningTok}[1]{\textcolor[rgb]{0.38,0.63,0.69}{\textbf{\textit{#1}}}}
\usepackage{longtable,booktabs,array}
\usepackage{calc} % for calculating minipage widths
% Correct order of tables after \paragraph or \subparagraph
\usepackage{etoolbox}
\makeatletter
\patchcmd\longtable{\par}{\if@noskipsec\mbox{}\fi\par}{}{}
\makeatother
% Allow footnotes in longtable head/foot
\IfFileExists{footnotehyper.sty}{\usepackage{footnotehyper}}{\usepackage{footnote}}
\makesavenoteenv{longtable}
\ifLuaTeX
\usepackage[bidi=basic]{babel}
\else
\usepackage[bidi=default]{babel}
\fi
\babelprovide[main,import]{french}
% get rid of language-specific shorthands (see #6817):
\let\LanguageShortHands\languageshorthands
\def\languageshorthands#1{}
\setlength{\emergencystretch}{3em} % prevent overfull lines
\providecommand{\tightlist}{%
  \setlength{\itemsep}{0pt}\setlength{\parskip}{0pt}}
% Ajouté au préambule de chaque document par generer.sh
\usepackage{xcolor}
\usepackage{newunicodechar}
\newfontfamily\symbolfont{DejaVuSans.ttf}
% Caractères absents de Latin Modern (flèches, comparateurs, cases) : repli sur DejaVu Sans
\newunicodechar{→}{{\symbolfont →}}
\newunicodechar{←}{{\symbolfont ←}}
\newunicodechar{⇄}{{\symbolfont ⇄}}
\newunicodechar{↔}{{\symbolfont ↔}}
\newunicodechar{≥}{{\symbolfont ≥}}
\newunicodechar{≤}{{\symbolfont ≤}}
\newunicodechar{≠}{{\symbolfont ≠}}
\newunicodechar{≈}{{\symbolfont ≈}}
\newunicodechar{✓}{{\symbolfont ✓}}
\newunicodechar{✗}{{\symbolfont ✗}}
\newunicodechar{⚠}{{\symbolfont ⚠}}
\newunicodechar{✅}{{\symbolfont\color[rgb]{0.1,0.55,0.2}✔}}
\newunicodechar{‑}{\mbox{-}}
% Retour à la ligne dans les blocs de code et les URL longues
\usepackage{fvextra}
\DefineVerbatimEnvironment{Highlighting}{Verbatim}{breaklines,commandchars=\\\{\}}
\usepackage{xurl}
\setlength{\emergencystretch}{3em}
\AtBeginEnvironment{longtable}{\small}
% Figures flottantes (pas de demi-pages vides), jamais plus hautes que 85 % de la page
\usepackage{float}
\floatplacement{figure}{htbp}
\usepackage{adjustbox}
\providecommand*\pandocbounded[1]{#1}% absent des documents sans image
\renewcommand*\pandocbounded[1]{\adjustbox{max width=\linewidth,max totalheight=0.85\textheight}{#1}}
\usepackage[font=small,labelfont=bf]{caption}
% Contenu sur une nouvelle page après la table des matières
\let\tdmorigine\tableofcontents
\renewcommand{\tableofcontents}{\tdmorigine\newpage}
\usepackage{fancyhdr}
\pagestyle{fancy}
\fancyhf{}
\fancyhead[L]{\small SmartRecruit — Qualité logicielle ISO/IEC 25010}
\fancyhead[R]{\small\thepage}
\renewcommand{\headrulewidth}{0.4pt}
\usepackage{bookmark}
\IfFileExists{xurl.sty}{\usepackage{xurl}}{} % add URL line breaks if available
\urlstyle{same}
\hypersetup{
  pdftitle={Fiche d\textquotesingle étude --- Membre 1 (lead) : présentation, logique métier et API},
  pdflang={fr},
  colorlinks=true,
  linkcolor={Maroon},
  filecolor={Maroon},
  citecolor={Blue},
  urlcolor={Blue},
  pdfcreator={LaTeX via pandoc}}

\title{Fiche d\textquotesingle étude --- Membre 1 (lead) : présentation,
logique métier et API}
\author{}
\date{}

\begin{document}
\maketitle

\renewcommand*\contentsname{Table des matières}
{
\hypersetup{linkcolor=}
\setcounter{tocdepth}{3}
\tableofcontents
}
\begin{quote}
Pour : \emph{(à compléter)} --- le membre qui connaît le mieux le
projet. Complète le \href{SOUTENANCE.pdf}{guide de soutenance}
(questions-réponses) et la \href{REPARTITION_TESTS.pdf}{répartition des
tests}.
\end{quote}

\section{0. Pourquoi cette partie revient au membre qui connaît le mieux
le
projet}\label{0-pourquoi-cette-partie-revient-au-membre-qui-connauxeet-le-mieux-le-projet}

\begin{itemize}
\tightlist
\item
  \textbf{Le lead présente tout} : sujet, architecture, choix
  techniques, démo, synthèse ISO 25010. Il faut une vue
  d\textquotesingle ensemble que les autres n\textquotesingle ont pas.
\item
  \textbf{La logique métier est la partie la plus difficile à apprendre
  vite} : 6 services, 9 règles métier, un workflow, 1 233 tests. Les
  autres blocs reposent surtout sur un outil et ses rapports (k6,
  Playwright, ZAP\ldots), qui s\textquotesingle apprennent en quelques
  jours ; la logique métier demande de comprendre le code.
\item
  \textbf{C\textquotesingle est ce que l\textquotesingle enseignant
  vérifie en premier} (§4 et §6.1 : « toutes les fonctionnalités
  implémentées et exactes »). Si une question dérape sur le code,
  c\textquotesingle est toi qui dois pouvoir répondre.
\item
  \textbf{Tu es le filet de sécurité} : quand un membre bloque, tu
  reprends (« en complément\ldots{} »). Tu dois donc connaître une
  phrase-clé de chaque autre partie (section 7).
\end{itemize}

\section{1. Ta mission en soutenance}\label{1-ta-mission-en-soutenance}

\begin{enumerate}
\def\labelenumi{\arabic{enumi}.}
\tightlist
\item
  Présenter le projet en ≈ 12 min (déroulé : \href{SOUTENANCE.pdf}{guide
  de soutenance §2}).
\item
  Faire la démo live de l\textquotesingle application et lancer les
  \textbf{2 couches de test en direct} (Playwright et k6).
\item
  Répondre sur ton bloc : \textbf{6.1 Aptitude fonctionnelle} et
  \textbf{6.7 Maintenabilité (couverture)} ; couches unitaires,
  intégration, API.
\item
  Répondre aux questions transverses (traçabilité, règle
  d\textquotesingle or, taille du groupe).
\end{enumerate}

\section{2. Ce que tu dois apprendre
(notions)}\label{2-ce-que-tu-dois-apprendre-notions}

\begin{longtable}[]{@{}
  >{\raggedright\arraybackslash}p{(\linewidth - 2\tabcolsep) * \real{0.3694}}
  >{\raggedright\arraybackslash}p{(\linewidth - 2\tabcolsep) * \real{0.6306}}@{}}
\toprule\noalign{}
\begin{minipage}[b]{\linewidth}\raggedright
Notion
\end{minipage} & \begin{minipage}[b]{\linewidth}\raggedright
Ce que tu dois savoir dire
\end{minipage} \\
\midrule\noalign{}
\endhead
\bottomrule\noalign{}
\endlastfoot
ISO/IEC 25010 & Modèle de qualité du produit logiciel, \textbf{8
caractéristiques} : aptitude fonctionnelle, efficacité de performance,
compatibilité, facilité d\textquotesingle utilisation, fiabilité,
sécurité, maintenabilité, portabilité. Pour chacune : une exigence
mesurable, un test, un outil, un résultat. \\
Traçabilité & Relier chaque exigence (EF, RG, EQ) à ses tests, son
outil, son résultat et sa preuve : c\textquotesingle est la
\href{MATRICE_TRACABILITE.pdf}{matrice}. \\
Pyramide des tests & Beaucoup de tests unitaires (rapides, ciblés),
moins d\textquotesingle intégration, peu d\textquotesingle E2E (lents,
coûteux). Chez nous : 1 233 unitaires/intégration, 155 API, 62 E2E. \\
Test unitaire & Une classe isolée (règle de scoring, validation), sans
base ni réseau. \\
Test d\textquotesingle intégration & Plusieurs couches ensemble :
requête HTTP simulée (MockMvc) → contrôleur → service → repository →
base H2. Les autres microservices sont remplacés par de \textbf{faux
serveurs HTTP}. \\
Test d\textquotesingle API (contrat) & Un vrai client HTTP contre la
stack Docker complète : codes HTTP et structure JSON. \\
Double de test (mock, stub, faux serveur) & Remplace une dépendance pour
isoler ce qu\textquotesingle on teste et rendre le test
reproductible. \\
Test paramétré & Un même test exécuté sur une liste de cas (ex. toutes
les paires de statuts pour RG‑03). \\
Codes HTTP & 400 requête invalide / transition interdite ; 401 non
authentifié ; 403 authentifié mais rôle insuffisant ; 404 introuvable ;
409 conflit avec l\textquotesingle état (doublon, offre non publiée) ;
422 données valides mais non traitables (profil incomplet) ; 429 trop de
requêtes. \\
Machine à états (workflow) & Des statuts et une table de transitions
autorisées ; tout le reste est refusé. \\
Couverture de code & \% de lignes (ou branches) exécutées par les tests
; mesure ce qui est \textbf{exécuté}, pas ce qui est
\textbf{vérifié}. \\
Architecture en couches & \texttt{controller} (HTTP) → \texttt{service}
(règles) → \texttt{repository} (accès base) → \texttt{model} (entités),
plus \texttt{dto}, \texttt{exception}, \texttt{client}. \\
Microservices & Un service par domaine, une base par service,
communication REST ; le gateway est la seule porte
d\textquotesingle entrée. \\
\end{longtable}

\section{3. Le projet vu de ta partie : fichiers à lire, dans
l\textquotesingle ordre}\label{3-le-projet-vu-de-ta-partie--fichiers-uxe0-lire-dans-lordre}

\begin{longtable}[]{@{}
  >{\raggedright\arraybackslash}p{(\linewidth - 4\tabcolsep) * \real{0.0155}}
  >{\raggedright\arraybackslash}p{(\linewidth - 4\tabcolsep) * \real{0.7674}}
  >{\raggedright\arraybackslash}p{(\linewidth - 4\tabcolsep) * \real{0.2171}}@{}}
\toprule\noalign{}
\begin{minipage}[b]{\linewidth}\raggedright
\#
\end{minipage} & \begin{minipage}[b]{\linewidth}\raggedright
Fichier
\end{minipage} & \begin{minipage}[b]{\linewidth}\raggedright
Ce que tu y cherches
\end{minipage} \\
\midrule\noalign{}
\endhead
\bottomrule\noalign{}
\endlastfoot
1 & Le cahier des charges (PDF de l\textquotesingle enseignant) puis
\href{RAPPORT_CONCEPTION.pdf}{RAPPORT\_CONCEPTION.md} & Le §4 du cahier
et comment chaque exigence est réalisée ; les 8 figures. \\
2 & \href{MATRICE_TRACABILITE.pdf}{MATRICE\_TRACABILITE.md} & Chaque
ligne EF/RG/EQ : savoir la lire à voix haute. \\
3 &
\texttt{backend/}\allowbreak{}\texttt{application-}\allowbreak{}\texttt{service/}\allowbreak{}\texttt{src/}\allowbreak{}\texttt{main/}\allowbreak{}\texttt{java/}\allowbreak{}\texttt{com/}\allowbreak{}\texttt{smartrecruit/}\allowbreak{}\texttt{application/}\allowbreak{}\texttt{model/}\allowbreak{}\texttt{Application}\allowbreak{}\texttt{Status.}\allowbreak{}\texttt{java}
& La table \texttt{ALLOWED\_TRANSITIONS} (RG‑03). \\
4 &
\texttt{…/}\allowbreak{}\texttt{application/}\allowbreak{}\texttt{service/}\allowbreak{}\texttt{Application}\allowbreak{}\texttt{Service.}\allowbreak{}\texttt{java}
& Postuler (RG‑04, RG‑05), changer de statut, historique, audit. \\
5 &
\texttt{…/}\allowbreak{}\texttt{application/}\allowbreak{}\texttt{controller/}\allowbreak{}\texttt{Application}\allowbreak{}\texttt{Controller.}\allowbreak{}\texttt{java}
et
\texttt{exception/}\allowbreak{}\texttt{Global}\allowbreak{}\texttt{Exception}\allowbreak{}\texttt{Handler.}\allowbreak{}\texttt{java}
& Comment une erreur métier devient un code HTTP (400/409/422). \\
6 &
\texttt{…/}\allowbreak{}\texttt{application/}\allowbreak{}\texttt{client/}\allowbreak{}
(\texttt{JobDirectoryClient},
\texttt{Candidate}\allowbreak{}\texttt{Readiness}\allowbreak{}\texttt{Client},
\texttt{AiMatchingClient}, \texttt{AuditClient}) & Les appels
inter-services ; l\textquotesingle IA en « best effort ». \\
7 &
\texttt{backend/}\allowbreak{}\texttt{application-}\allowbreak{}\texttt{service/}\allowbreak{}\texttt{src/}\allowbreak{}\texttt{test/}\allowbreak{}\texttt{java/}\allowbreak{}\texttt{…/}\allowbreak{}\texttt{controller/}\allowbreak{}\texttt{Status}\allowbreak{}\texttt{Change}\allowbreak{}\texttt{Integration}\allowbreak{}\texttt{Test.}\allowbreak{}\texttt{java}
et
\texttt{Apply}\allowbreak{}\texttt{Integration}\allowbreak{}\texttt{Test.}\allowbreak{}\texttt{java}
& Structure d\textquotesingle un test d\textquotesingle intégration :
préparation, appel MockMvc, assertions. \\
8 &
\texttt{backend/}\allowbreak{}\texttt{auth-}\allowbreak{}\texttt{service/}\allowbreak{}\texttt{src/}\allowbreak{}\texttt{test/}\allowbreak{}\texttt{java/}\allowbreak{}\texttt{…/}\allowbreak{}\texttt{controller/}\allowbreak{}\texttt{Login}\allowbreak{}\texttt{Integration}\allowbreak{}\texttt{Test.}\allowbreak{}\texttt{java}
& RG‑06 (verrouillage après 5 échecs). \\
9 &
\texttt{ai-}\allowbreak{}\texttt{service/}\allowbreak{}\texttt{app/}\allowbreak{}\texttt{services/}\allowbreak{}\texttt{scoring.}\allowbreak{}\texttt{py},
\texttt{matching\_service.py}, \texttt{rules\_match.py} +
\texttt{ai-}\allowbreak{}\texttt{service/}\allowbreak{}\texttt{tests/}\allowbreak{}\texttt{test\_}\allowbreak{}\texttt{scoring.}\allowbreak{}\texttt{py}
& Comment le score de compatibilité est calculé et testé. \\
10 &
\texttt{backend/}\allowbreak{}\texttt{e2e/}\allowbreak{}\texttt{e2e\_}\allowbreak{}\texttt{full.}\allowbreak{}\texttt{py}
& Les 155 vérifications API, en passant par le gateway et nginx. \\
11 & \texttt{quality/run-all.sh} & Les 10 étapes de la campagne
(install, unit, api, ui, perf, reliability, sast, sca, dast,
maintainability). \\
\end{longtable}

\section{4. Commandes à savoir lancer (et
lire)}\label{4-commandes-uxe0-savoir-lancer-et-lire}

\begin{Shaded}
\begin{Highlighting}[]
\FunctionTok{bash}\NormalTok{ quality/run{-}all.sh unit}
\end{Highlighting}
\end{Shaded}

Tests unitaires + intégration des 6 services Java, du service IA et du
frontend, avec couverture. Lire : le nombre de tests et
\texttt{BUILD\ SUCCESS} / \texttt{passed} ; rapports dans
\texttt{quality/}\allowbreak{}\texttt{reports/}\allowbreak{}\texttt{unit/}\allowbreak{}
et
\texttt{quality/}\allowbreak{}\texttt{reports/}\allowbreak{}\texttt{coverage/}\allowbreak{}.

\begin{Shaded}
\begin{Highlighting}[]
\BuiltInTok{cd}\NormalTok{ ai{-}service }\KeywordTok{\&\&} \ExtensionTok{pytest} \AttributeTok{{-}q}
\end{Highlighting}
\end{Shaded}

Version rapide (quelques secondes, 137 tests) : ton \textbf{plan B} si
une démo échoue.

\begin{Shaded}
\begin{Highlighting}[]
\FunctionTok{bash}\NormalTok{ quality/run{-}all.sh api}
\end{Highlighting}
\end{Shaded}

Application démarrée. Lire la dernière ligne de
\texttt{quality/}\allowbreak{}\texttt{reports/}\allowbreak{}\texttt{api/}\allowbreak{}\texttt{e2e-}\allowbreak{}\texttt{console.}\allowbreak{}\texttt{txt}
: 155 vérifications OK.

Démo de la soutenance (tu les lances, les membres 2 et 3 les expliquent
si on leur demande) :

\begin{Shaded}
\begin{Highlighting}[]
\BuiltInTok{cd}\NormalTok{ quality/e2e{-}ui }\KeywordTok{\&\&} \ExtensionTok{npx}\NormalTok{ playwright test }\AttributeTok{{-}{-}project}\OperatorTok{=}\NormalTok{chromium }\AttributeTok{{-}{-}headed}
\end{Highlighting}
\end{Shaded}

\begin{Shaded}
\begin{Highlighting}[]
\ExtensionTok{docker}\NormalTok{ run }\AttributeTok{{-}{-}rm} \AttributeTok{{-}v} \StringTok{"}\VariableTok{$(}\BuiltInTok{pwd} \AttributeTok{{-}W} \DecValTok{2}\OperatorTok{\textgreater{}}\NormalTok{/dev/null }\KeywordTok{||} \BuiltInTok{pwd}\VariableTok{)}\StringTok{/quality/perf:/scripts:ro"} \AttributeTok{{-}v} \StringTok{"}\VariableTok{$(}\BuiltInTok{pwd} \AttributeTok{{-}W} \DecValTok{2}\OperatorTok{\textgreater{}}\NormalTok{/dev/null }\KeywordTok{||} \BuiltInTok{pwd}\VariableTok{)}\StringTok{/quality/reports/perf:/reports"}\NormalTok{ grafana/k6:1.0.0 run /scripts/load{-}test.js}
\end{Highlighting}
\end{Shaded}

\section{5. Chiffres à connaître par
cœur}\label{5-chiffres-uxe0-connauxeetre-par-cux153ur}

\begin{longtable}[]{@{}
  >{\raggedright\arraybackslash}p{(\linewidth - 2\tabcolsep) * \real{0.2373}}
  >{\raggedright\arraybackslash}p{(\linewidth - 2\tabcolsep) * \real{0.7627}}@{}}
\toprule\noalign{}
\begin{minipage}[b]{\linewidth}\raggedright
Sujet
\end{minipage} & \begin{minipage}[b]{\linewidth}\raggedright
Valeur
\end{minipage} \\
\midrule\noalign{}
\endhead
\bottomrule\noalign{}
\endlastfoot
Tests unitaires + intégration & \textbf{1 233} = 693 JUnit (6 services)
+ 137 pytest (IA) + 403 Vitest (front) \\
Tests API & \textbf{155/155} \\
Couverture des lignes & Java \textbf{95,7 \%} (91,4 à 97,7 \% selon le
service), IA \textbf{97,2 \%}, frontend \textbf{92,1 \%} \\
Règles métier & \textbf{9} (RG‑01 à RG‑09), au moins un test nommé par
règle \\
Volume pour la recherche & \textbf{3 068} offres publiées \\
Campagne finale & 11/11/2026, toutes les étapes réussies en un passage
(\texttt{quality/}\allowbreak{}\texttt{reports/}\allowbreak{}\texttt{campagne-}\allowbreak{}\texttt{finale-}\allowbreak{}\texttt{console.}\allowbreak{}\texttt{txt}) \\
Les autres blocs (pour la synthèse) & E2E 62/62 sur 4 navigateurs ; p95
349,6 ms \textless{} 500 ms ; fiabilité 8/8 après 2 défauts corrigés ;
CVE 179 → 0 ; ZAP 0 échec ; CI 20/20 jobs \\
\end{longtable}

\section{6. Preuves à savoir ouvrir}\label{6-preuves-uxe0-savoir-ouvrir}

\begin{itemize}
\tightlist
\item
  \texttt{quality/}\allowbreak{}\texttt{reports/}\allowbreak{}\texttt{coverage/}\allowbreak{}\texttt{job-}\allowbreak{}\texttt{service/}\allowbreak{}\texttt{index.}\allowbreak{}\texttt{html}
  (rapport JaCoCo : montrer une classe verte/rouge).
\item
  \texttt{quality/}\allowbreak{}\texttt{reports/}\allowbreak{}\texttt{unit/}\allowbreak{}\texttt{\textless{}service\textgreater{}/}\allowbreak{}
  (résultats JUnit) et
  \texttt{quality/}\allowbreak{}\texttt{reports/}\allowbreak{}\texttt{unit/}\allowbreak{}\texttt{ai-}\allowbreak{}\texttt{service-}\allowbreak{}\texttt{junit.}\allowbreak{}\texttt{xml}.
\item
  \texttt{quality/}\allowbreak{}\texttt{reports/}\allowbreak{}\texttt{api/}\allowbreak{}\texttt{e2e-}\allowbreak{}\texttt{console.}\allowbreak{}\texttt{txt}.
\item
  \href{MATRICE_TRACABILITE.pdf}{MATRICE\_TRACABILITE.md} et
  \href{RAPPORT_CONFORMITE_ISO25010.pdf}{RAPPORT\_CONFORMITE\_ISO25010.md}
  : savoir aller directement à une ligne demandée.
\end{itemize}

\section{7. Le minimum à savoir sur les autres parties (filet de
sécurité)}\label{7-le-minimum-uxe0-savoir-sur-les-autres-parties-filet-de-suxe9curituxe9}

\begin{longtable}[]{@{}
  >{\raggedright\arraybackslash}p{(\linewidth - 2\tabcolsep) * \real{0.3966}}
  >{\raggedright\arraybackslash}p{(\linewidth - 2\tabcolsep) * \real{0.6034}}@{}}
\toprule\noalign{}
\begin{minipage}[b]{\linewidth}\raggedright
Partie
\end{minipage} & \begin{minipage}[b]{\linewidth}\raggedright
Une phrase
\end{minipage} \\
\midrule\noalign{}
\endhead
\bottomrule\noalign{}
\endlastfoot
Membre 2 --- Interface & Playwright rejoue 13 parcours réels sur
Chromium, Edge, WebKit et Chrome mobile ; axe-core : 0 violation WCAG
2.1 A/AA sur 6 pages ; Firefox documenté non exécutable (règle
d\textquotesingle or). \\
Membre 3 --- Performance, fiabilité, portabilité & k6 : 20 utilisateurs
pendant 1 min, p95 349,6 ms ; 8 pannes réelles de conteneurs, 2 défauts
trouvés et corrigés ; tout s\textquotesingle installe avec
\texttt{docker\ compose\ up\ -}\allowbreak{}\texttt{d\ -}\allowbreak{}\texttt{-}\allowbreak{}\texttt{build},
CI 20/20 jobs. \\
Membre 4 --- Sécurité, qualité de code & SAST 0 (Bandit, Semgrep), SCA
179 → 0 CVE (Trivy), ZAP 0 échec ; JWT + RBAC + en-têtes signés HMAC ;
CCN moyen 2,2, lint à 0. \\
\end{longtable}

\section{8. Questions pièges (au-delà du guide de
soutenance)}\label{8-questions-piuxe8ges-au-deluxe0-du-guide-de-soutenance}

\begin{longtable}[]{@{}
  >{\raggedright\arraybackslash}p{(\linewidth - 2\tabcolsep) * \real{0.5000}}
  >{\raggedright\arraybackslash}p{(\linewidth - 2\tabcolsep) * \real{0.5000}}@{}}
\toprule\noalign{}
\begin{minipage}[b]{\linewidth}\raggedright
Question
\end{minipage} & \begin{minipage}[b]{\linewidth}\raggedright
Réponse
\end{minipage} \\
\midrule\noalign{}
\endhead
\bottomrule\noalign{}
\endlastfoot
Pourquoi des microservices pour un projet étudiant ? & Un service par
domaine (comptes, profils, offres, candidatures, statistiques, IA) :
responsabilités séparées et testables isolément, et une panne
d\textquotesingle un service n\textquotesingle arrête pas les autres
(prouvé par les tests de fiabilité). Le prix : de la complexité,
compensée par Docker Compose. \\
Comment garantir la cohérence sans clés étrangères entre bases ? & Clés
étrangères logiques + vérifications par appel inter-service avant
d\textquotesingle écrire (ex. l\textquotesingle offre existe et est
\texttt{PUBLISHED} avant d\textquotesingle accepter une candidature). \\
Différence entre 409 et 422 dans votre API ? & 409 : conflit avec
l\textquotesingle état actuel (doublon, offre non publiée). 422 : la
requête est correcte mais le candidat ne remplit pas une condition
(profil incomplet, pas de CV). \\
Un test qui passe toujours, comment savoir qu\textquotesingle il teste
quelque chose ? & Chaque règle a une assertion explicite sur le résultat
\textbf{et} sur l\textquotesingle état (ex. RG‑03 : 400 \textbf{et}
statut inchangé). On a vu des tests échouer avant correction (fiabilité
5/8 au premier passage). \\
Où est l\textquotesingle exigence « CRUD complet sur 4 entités » ? &
Rapport de conception §3 : 6 entités avec les 4 opérations ; tests
\texttt{Job}\allowbreak{}\texttt{Delete}\allowbreak{}\texttt{Integration}\allowbreak{}\texttt{Test},
\texttt{Notes}\allowbreak{}\texttt{Integration}\allowbreak{}\texttt{Test}\ldots{}
(matrice EF‑02). \\
Que se passe-t-il si l\textquotesingle IA est en panne pendant une
candidature ? & La candidature est enregistrée (201), le score reste
\texttt{null} : l\textquotesingle IA est appelée en « best effort » et
de façon asynchrone (figure 7 du rapport de conception). \\
Votre IA est-elle testée alors qu\textquotesingle elle dépend
d\textquotesingle un LLM externe ? & Le scoring par règles est
déterministe et testé (pytest, 97,2 \% de couverture) ;
l\textquotesingle appel au LLM externe est testé avec un client simulé
(\texttt{test\_llm\_client.py}). \\
\end{longtable}

\section{9. Plan
d\textquotesingle apprentissage}\label{9-plan-dapprentissage}

\begin{longtable}[]{@{}
  >{\raggedright\arraybackslash}p{(\linewidth - 4\tabcolsep) * \real{0.1316}}
  >{\raggedright\arraybackslash}p{(\linewidth - 4\tabcolsep) * \real{0.1316}}
  >{\raggedright\arraybackslash}p{(\linewidth - 4\tabcolsep) * \real{0.7368}}@{}}
\toprule\noalign{}
\begin{minipage}[b]{\linewidth}\raggedright
Étape
\end{minipage} & \begin{minipage}[b]{\linewidth}\raggedright
Durée
\end{minipage} & \begin{minipage}[b]{\linewidth}\raggedright
Objectif
\end{minipage} \\
\midrule\noalign{}
\endhead
\bottomrule\noalign{}
\endlastfoot
1 & 2 h & Relire le cahier des charges, le rapport de conception et la
matrice : savoir raconter le projet sans notes. \\
2 & 3 h & Lire les fichiers 3 à 8 de la section 3 ; refaire à voix haute
le parcours « postuler » (figure 7). \\
3 & 1 h & Lancer \texttt{unit} et \texttt{api}, ouvrir chaque preuve de
la section 6. \\
4 & 2 h & Répéter la présentation de 12 min avec la démo et les 2
commandes live, chronomètre en main. \\
5 & 1 h & Faire interroger chaque membre par un autre avec le guide de
soutenance ; noter ce qui bloque. \\
\end{longtable}

\section{10. Auto-évaluation (tout doit être coché la
veille)}\label{10-auto-uxe9valuation-tout-doit-uxeatre-cochuxe9-la-veille}

\begin{itemize}
\tightlist
\item[$\square$]
  Je cite les 8 caractéristiques ISO 25010 et une preuve pour chacune.
\item[$\square$]
  Je dessine l\textquotesingle architecture de mémoire (gateway, 6
  services Spring, IA FastAPI, une base par service).
\item[$\square$]
  J\textquotesingle explique RG‑03 et RG‑04 en montrant le code
  \textbf{et} le test.
\item[$\square$]
  Je justifie 400 / 401 / 403 / 409 / 422 / 429 avec un exemple de notre
  API.
\item[$\square$]
  Je lance \texttt{unit}, \texttt{api}, Playwright et k6 sans aide, et
  je sais quoi lire dans leur sortie.
\item[$\square$]
  Je retrouve en moins de 30 s n\textquotesingle importe quelle ligne de
  la matrice et sa preuve.
\item[$\square$]
  Je tiens la présentation en 12 min.
\end{itemize}

\end{document}
